\documentclass[main.tex]{subfiles}


\begin{document}

% \AddToShipoutPictureFG{%
% \begin{tikzpicture}[overlay,remember picture]
% \path (current page.south west) -- (current page.north east)
% node[midway,scale=1,color=gray,sloped,opacity=0.4] {\textnormal{Кравченко Игорь Игоревич\hspace{1cm} +79010144910\hspace{1cm} super.antig@yandex.ru}};
% \end{tikzpicture}
% }


% %from pagenumber to up
% \phantomsection 
% \hypertarget{toc}{}

 

 

 
   

\section{Основные формулы молекулярной физики}


В молекулярной физике рассматриваются \emph{макроскопические} тела~--- то есть тела, состоящие из огромного числа частиц\footnote{Обычно речь идет о молекулах или атомах, что предполагается далее.}: например, даже в одном грамме воды содержится порядка~$10^{22}$ молекул. Для удобства описания таких тел с учетом их внутреннего строения вводят следующие величины.
\begin{itemize}
    \item \textbf{Число Авогадро} $\left(N_\text{А}\ \left[\dfrac{1}{\text{моль}}\right]\right)$~--- это число частиц, содержащееся в такой порции вещества, \strut которой удобно пользоваться на практике (эта порция носит название <<моль>>). Моль любого вещества содержит одно и то же число частиц, которое в настоящее время принимают приближенно равным  $N_\text{А}\approx6\cdot10^{23}$~моль$^{-1}$.
    
    \item \textbf{Количество вещества} ($\nu$ [моль])~--- это количество порций по $N_\text{А}$ частиц, образующих данное тело. 

    \item \textbf{Молярная масса} $\left(M\ \left[\dfrac{\text{кг\vphantom{1}}}{\text{моль}}\right]\right)$~--- это масса порции из $N_\text{А}$ частиц данного вещества (или, как говорят, \strut масса одного моля этого вещества). Молярные массы химических элементов можно узнать из таблицы Менделеева: следует взять атомную массу (число возле номера элемента) из ячейки данного элемента и умножить на~$10^{-3}$~--- получится значение в $\text{кг}/\text{моль}$. Например, в ячейке железа у его порядкового номера стоит число~56 (округленно); это значит, что молярная масса железа (одноатомное вещество) равна~$56\cdot10^{-3}$~$\text{кг}/\text{моль}$.

    Для определения молярной массы вещества, молекулы которого состоят из нескольких атомов, нужно суммировать молярные массы. Например, молярная масса воды~H$_2$O равна $1\cdot10^{-3}\cdot2+16\cdot10^{-3}=18\cdot10^{-3}$~$\text{кг}/\text{моль}$. 
    \end{itemize}

\textbf{Число частиц} в теле с учетом определения количества вещества можно вычислять по формуле:
\begin{equation}
    N=\nu N_\text{А}.
\end{equation}

\textbf{Масса} вещества с учетом определения молярной массы рассчитывается так:
\begin{equation}
    m=\nu M,
\end{equation}
или, обозначая массу частицы через~$m_0$:
\begin{equation}
    m=Nm_0.
\end{equation}

\textbf{Плотность}, как известно из механики, находится по формуле:
\begin{equation}
    \rho=\frac{m}{V}.
\end{equation}

\textbf{Концентрация} $\left(n\ \left[\dfrac{1}{\text{м}^3}\right]\right)$~--- это характеристика тела, показывающая число частиц в единице \strut объема:
\begin{equation}
    n=\dfrac{N}{V}.
\end{equation}


































%%%%%%%%%%%%%
% далее посчитался ненужным такой листок, но инфа может пригодится далее при упоминании идеального газа

% \textbf{Идеальный газ}~--- это \emph{физическая модель}\footnote{Физическая модель~--- это упрощенный аналог физической системы (процесса), сохраняющий ее главные черты.}, использующаяся для описания разреженных газов, когда расстояния между их частицами\footnote{Молекулами или атомами.} намного больше размеров самих частиц. Эта модель предполагает следующие допущения.
% \begin{enumerate}
%     \item Частицы газа считаются материальными точками (рис.\,\ref{fig:p68}).

%     \begin{figure}[H]
%     \centering

%     \begin{tikzpicture}[font=\small,            line cap=round,                             >={Stealth[inset=0pt,round,length=3mm, angle'=20]},vec/.style={>={Stealth[inset=0pt,round,length=3.5mm, angle'=20]}}]


% %gas
% \draw (1,0) rectangle++ (2,2);

% % \foreach \i in {1,...,100}
% %   \fill[Green,shift={(1.03,0.03)}] ({rnd*1.94cm}, {rnd*1.94cm}) circle (.02);


%     \end{tikzpicture}
% \caption{Размеры частиц пренебрежимо малы}
% \label{fig:p68}
% \end{figure}

%     \item Частицы не взаимодействуют друг с другом на расстоянии.

% \tikzfading[name=fade right,
% left color=transparent!50,
% right color=transparent!100]

%     \begin{figure}[H]
%     \centering

%     \begin{tikzpicture}[font=\small,            line cap=round,                             >={Stealth[inset=0pt,round,length=3mm, angle'=20]},vec/.style={>={Stealth[inset=0pt,round,length=3.5mm, angle'=20]}}]


% \def\num{2}

% \foreach \i in {1,...,\num}
%   \path[Green] (rnd*360:rnd*2cm) circle (.05) coordinate (m\i);

% \foreach \i in {1,...,\num}{
%   \def\ang{rnd*360}
%   \draw[Green,path fading=fade right,fading angle=\ang,line width=.05cm] (m\i) --++ (rnd*\ang:1cm);
% };

% \foreach \i in {1,...,\num}
%   \fill[Green] (m\i) circle (.05);

% % \draw[Green,path fading=fade right,fading angle=90,line width=.05cm] (0,0) --++ (90:1cm);


% % \node at (m1) {1};

%     \end{tikzpicture}
% \caption{Размеры частиц пренебрежимо малы}
% \label{fig:p69}
% \end{figure}
    
%     \item Столкновения частиц считают абсолютно упругими.
% \end{enumerate}

















 
\end{document}